

\vspace{-2.0em}
\noindent\uppercase{Der vierte Band der Reihe VIII}
\vspace{0.5em}
\par\noindent
Für den hier nun vorliegenden vierten Band der naturwissenschaftlichen, medizinischen und technischen Schriften von Leibniz gilt, was vor vier Jahren  für den dritten Band der Reihe festzuhalten war: Bis auf die Nachtragsstücke vereint er Schriften, die Leibniz nach seiner Rückkehr aus Frankreich verfasst hat. Sie stammen aus einem Zeitraum, der in diesem Band von 1676 bis ins letzte Lebensjahrzehnt reicht. Eine Datierung der vielfach undatierten Schriften ist dadurch schwierig. Auch inhaltlich wie editorisch stellt dieser Band wieder hohe Anforderungen, weil er sich wie sein Vorgänger auf Stücke aus den Bereichen der Mechanik beschränkt. Bei diesen Themen begnügt sich Leibniz nicht damit, sie rein sprachlich auf Lateinisch, Französisch oder Deutsch zu behandeln, sondern sucht sie mathematisch, gestützt auf Diagramme zu durchdringen.%
\\ \indent
Wie sein Vorgänger folgt der vierte Band der Reihe dem Prinzip der Modularisierung: Das größte Modul bildet hierbei die Mechanik, das in weiteren Bänden bearbeitet werden wird. Eine chronologische Ordnung der edierten Handschriften ist damit nur in den Untermodulen der Mechanik möglich. Sie erfolgt hier jeweils für die drei Abschnitte des Bandes, die aus den Stücken zum Stoß, zur Statik und den Stücken zur Akustik, Elastizität und Festigkeit (AEF) gebildet werden. Die Stücke zum Stoß finden in diesem Band eine geplante und angekündigte Fortsetzung zum gleichnamigen Abschnitt im letzten Band. Weitgehend gelingt diese Fortsetzung auch chronologisch, so dass der dritte Band die frühe Beschäftigung mit den Phänomenen des Stoßes dokumentiert, allerdings mit einer weiten zeitlichen Überlappung für die 1680er Jahre zu den Stoß-Stücken dieses Bandes. An den Stücken zur AEF bestätigt sich hingegen, dass das zu edierende, überwiegend unveröffentlichte und undatierte Material inhaltlich oft nur nach und nach richtig eingeordnet werden kann. Die AEF-Stücke in diesem Band komplettieren den gleichnamigen Abschnitt des vorangegangenen, stammen aber wie diese aus dem gesamten Zeitraum nach Paris. Zu hoffen bleibt, dass es einmal gelingen wird, diese vom Bearbeitungsfortschritt abhängige und in einzelne Bände getrennte Präsentation der Forschungsergebnisse auf einer gemeinsamen digitalen Basis zu vereinen, um die Stücke der Edition übergreifend nach Inhalt oder Entstehungszeit nutzen zu können.%
\newpage
\vspace{6.0em}
\noindent\uppercase{Stücke des Bandes im Überblick}
\par\vspace{1.0em}
\noindent
I. Akustik, Elastizität, Festigkeit (N.\,\ref{60250}--N.\,\ref{37109+60209})%
\par
\vspace{0.25em}
\noindent
Die unter dieser Überschrift vereinten Stücke zeichnen sich durch eine große thematische Breite aus, die es nicht erlaubt, sie eindeutig nur einer Betrachtung zuzuschreiben und ganz zu trennen. 
Wenn die Stücke hier nachfolgend in Gruppen eingeordnet präsentiert werden, geschieht dies unter Vorbehalt und im Bewusstsein, ihnen nicht immer ganz gerecht damit zu werden.%
\par
\vspace{1.0em}
\noindent
I.\;1. Materietheorie und Kohäsion%
\par \noindent
Die frühesten Stücke dieser ersten Gruppe befassen sich mit grundlegenden Fragen der Materie und des Zusammenhalts von Körpern. In einer kurzen Notiz argumentiert Leibniz, dass der Zusammenhalt der Körper durch Bewegung entstehe, wobei Bläschen als Grundbausteine der Materie fungierten (N.\,\ref{60250}). Diese durch die Bewegung des Äthers gleichsam belebten Bläschen bilden nach seiner Theorie die Grundlage für die Konsistenz und Vielfalt der Körper. Die Bedeutung dieser Hypothese bestätigt sich an einem Zitat aus Erhard Weigels\indexN{\textls{Weigel}, Erhard 1625\textendash1699} Dissertation, wonach die meisten Phänomene der Körper laut Leibniz durch Bläschen erklärt werden könnten (N.\,\ref{38575}).
\\ \indent
In zwei zusammengehörigen Aufzeichnungen aus den Jahren 1680 bis 1687 greift Leibniz seine Materietheorie wieder auf und setzt sich kritisch mit cartesischen Positionen auseinander. Er widerlegt zunächst die These, dass Ruhe die Ursache des Zusammenhalts der Körper sei, und hält stattdessen die Bewegung, insbesondere die der umgebenden Fluide, für die eigentliche Ursache der Kohäsion (N.\,\ref{60349-1}). Seine Argumentation stützt er auf Beispiele aus der Hydraulik und Pneumatik. In einer weiteren Aufzeichnung wendet er sich gegen Descartes'\indexN{\textls{Descartes} (Cartesius, des Cartes), René 1596\textendash1650} Elementenlehre  (N.\,\ref{60349-2}). Leibniz bezweifelt, dass Licht das feinste Element sei, und vertritt dagegen eine Theorie, wonach das Licht eine spezifische Flüssigkeit sei, deren Bewegung man sich ähnlich wie die des Schalls vorstellen könne.
\\ \indent
Die Auseinandersetzung mit grundlegenden Materieeigenschaften setzt sich in seinen Überlegungen zur Dichte fort; den Zusammenhang zwischen Dichte und Volumen betrachtet er als ein mathematisches Verhältnis, das der Kompression der Materie entspricht (N.\,\ref{58656-A}, N.\,\ref{58216}, N.\,\ref{58258}). In einer kurzen, aber grundlegenden Notiz unterscheidet er dynamische und plastische Naturgesetze, wobei er die universelle Elastizität der Materie als Voraussetzung für die dynamischen Gesetze postuliert (N.\,\ref{39746}). In einer weiteren Aufzeichnung erklärt Leibniz die Transparenz von Flüssigkeiten damit, dass sich durch ihre regelmäßig angeordneten Poren die Teilchen schnell bewegen können (N.\,\ref{60255}). Besondere Aufmerksamkeit schenkt Leibniz dem Verhalten von Wasser bei Kälte, wobei er die von Descartes\indexN{\textls{Descartes} (Cartesius, des Cartes), René 1596\textendash1650} vorgebrachte Erklärung zurückweist (N.\,\ref{60248}).
\par 
\vspace{1em}
\noindent
I.\;2. Akustik und Schallphänomene%
\par
\noindent
Akustische Phänomene im engeren Sinne sind Gegenstand von drei Stücken, wobei der Donner und das Echo im Mittelpunkt stehen. In einem auf Deutsch verfassten Text beschreibt Leibniz das Phänomen des Donners als eine Art Echo in den Wolken und untermauert diese Erklärung durch Beobachtungen von Schallphänomenen, insbesondere durch die Beschreibung eines bemerkenswerten Echos bei Helmstedt (N.\,\ref{56218}). In einer späteren Notiz ergänzt er diese Auffassung durch Beobachtungen, die den Zusammenhang zwischen Donner und Regen sowie zwischen Blitzstärke und Donnerlautstärke bestätigen (N.\,\ref{60249}). %
Im praktisch-technischen Bereich befasst sich Leibniz mit der Konstruktion von Musikinstrumenten. Eine Notiz beschreibt die optimale Konstruktion eines Glockenspiels, wobei er das Hebelgesetz anwendet, um eine ausgewogene Spielbarkeit aller Glocken unabhängig von ihrer Größe zu gewährleisten (N.\,\ref{57037-3}).%
\par
\vspace{1em}
\noindent
I.\;3. Barometrie und Thermometrie
\par
\noindent
Eine Reihe von Stücken beschäftigt sich damit, wie Mess- und Beobachtungsinstrumente verbessert werden können. Leibniz stellt Berechnungen an, um die Volumenverhältnisse in Barometern zu bestimmen (N.\,\ref{57960}). Ausgehend von Überlegungen bei  Christiaan Huygens\indexN{\textls{Huygens} (Hugenius, Ugenius, Hugens, Huguens), Christiaan 1629\textendash1695} und Louis Hubin\indexN{\textls{Hubin} (Hubinus), Louis (Ludion de) fl. 1673-1683} entwirft Leibniz ein tragbares kombiniertes Barometer-Thermometer (N.\,\ref{57037-1}). Auf diese Konzeption greift er in zwei aufeinanderfolgenden Aufzeichnungen zurück, analysiert die geometrischen und physikalischen Beziehungen zwischen Quecksilbersäule und Luftdruck und zieht daraus Erkenntnisse für die technische Umsetzung einer Konstruktion (N.\,\ref{57038}, N.\,\ref{57039}). %
Ein solche schlägt er sehr viel später für ein geschlossenes Thermometer vor, das frei von Einflüssen des Luftdrucks nur noch den  
temperaturbedingten Veränderungen der eingeschlossenen Luft unterliegt, die Leibniz mathematisch erfasst, um damit eine exakte Skalierung für das Instrument zu gewinnen (N.\,\ref{41212}, N.\,\ref{60326}). Mit Blick auf diese in der Elastizität begründeten Faktoren macht Leibniz Einwände gegen bzw. Verbesserungsvorschläge für eine von Guillaume Amonton\indexN{\textls{Amontons}, Guillaume 1663\textendash1705} vorgeschlagene Barometer-Konstruktion, die auf mehreren durch luftgefüllte Röhren verbundenen Quecksilbersäulen beruht (N.\,\ref{53451}).
Das barometrische Paradoxon bildet den Gegenstand einer Abfolge von Stücken. Wie andere Gelehrte seiner Zeit und teils in Auseinandersetzung mit ihnen -- Bernardino Ramazzini\indexN{\textls{Ramazzini} (Ramazzinus), Bernardino 1633\textendash1714}, Günther Christoph Schelhammer\indexN{\textls{Schelhammer} (Schelhammerus), Günther Christoph 1649\textendash1716}, Rudolf Jakob Camerarius\indexN{\textls{Camerarius}, Rudolf Jakob 1665\textendash1721} -- sucht Leibniz physikalisch zu begründen, warum das Gewicht der Luft (Luftdruck) bei schlechtem Wetter sinkt, bei gutem steigt. Für seine eigene Erklärung stützt er sich auf eine hydrostatische Annahme, derzufolge ein Körper mit seinem Gewicht eine Flüssigkeit nicht belaste, solange er in ihr versinke (N.\,\ref{53442}, N.\,\ref{60267}, N.\,\ref{36805}); in einer späteren Veröffentlichung, die 1715 in den von der Akademie der Naturforscher herausgegebenen \textit{Ephemeriden} erscheint, verteidigt Leibniz seine Klärung des Paradoxons gegen Einwände (N.\,\ref{54958}).%
\par 
\vspace{1em}
\noindent
I.\;4. Verhalten gasförmiger, fester und flüssiger Körper%
\par
\noindent
Eine Reihe von Stücken hat das Verhalten von Luft zum Gegenstand. In einer umfangreichen Aufzeichnung befasst sich Leibniz mit verschiedenen Experimenten im Vakuum und zieht daraus Rückschlüsse auf die elastischen Eigenschaften der Luft (N.\,\ref{60212}). Er diskutiert das Verhältnis zwischen Druck, Volumen und elastischer Kraft und leitet mathematische Gesetzmäßigkeiten ab. Zwei zusammengehörige Aufzeichnungen liefern Ansätze, um die elastische Kraft komprimierter Luft mathematisch zu bestimmen, wobei Leibniz zwischen "toter" und "lebendiger" Kraft unterscheidet (N.\,\ref{41212}, N.\,\ref{41213}).%
\par
Seine Beschäftigung mit den Eigenschaften der Luft ergänzt Leibniz durch weitere Experimente. Er stellt einen detaillierten Versuchsplan auf, um den Einfluss von Flammen auf die Elastizität der Luft zu beobachten (N.\,\ref{55743+55744+60656}). In Notizen zu einer fremdhändigen Aufzeichnung hält er fest, wie Wärme sich quantitativ auf die Ausdehnung von Luft auswirkt (N.\,\ref{36650}), und in einer weiteren, dass Luft aus zwei verschiedenen Materiearten bestehen könnte (N.\,\ref{56214}). Besondere Aufmerksamkeit kommt der Kraftübertragung auf Luft zu, wobei Leibniz argumentiert, dass Flügel aufgrund der Teilbarkeit der Luft effektiver sind als unmittelbare Luftströme (N.\,\ref{55805}). In einem technisch orientierten Text zieht er noch die Möglichkeit eines perpetuum mobile in Erwägung, das die Kraft des Windes nutzt (N.\,\ref{56215}). Diese Überlegungen führen ihn auch zu einem Konzept für eine künstliche fliegende Taube, die komprimierte Luft als Antrieb nutzen soll (N.\,\ref{55774}).%
\\ \indent
Ein eigenes Thema ist das Phänomen des Verdampfens. In zwei eng zusammenhängenden Stücken befasst Leibniz sich mit der Theorie dieses Phänomens. Gegen die von Etienne Chauvin\indexN{\textls{Chauvin} (Chauvinus), Etienne 1640--1725} vorgebrachte Erklärung bezieht Leibniz kritisch Stellung (N.\,\ref{60210}); er selbst führt das Verdampfen von Wasser auf Hohlräume zurück, die sich in den Wasser-Bläschen infolge der Wärmezufuhr gebildet haben und sie in der Luft aufsteigen lassen (N.\,\ref{37109+60209}). Praktische Anwendungen dieses Phänomens finden sich in seinen Kommentaren zu Denis Papins\indexN{\textls{Papin} (Papinus), Denis 1647\textendash1713} Abhandlung über die Gewinnung mechanischer Kräfte durch Dampf (N.\,\ref{55804}).%
\\ \indent
Die Mechanik gespannter Fäden und elastischer Körper ist ein weiteres Thema für Leibniz. Er  beweist, dass eine Saite mit Eigengewicht nicht perfekt gespannt werden kann (N.\,\ref{60198}). Um die Fadenspannung zu bestimmen, stellt er eine infinitesimale Berechnung an (N.\,\ref{57578}). Er definiert das Spannungszentrum an einer Geraden und beweist dessen Analogie zum Schwerpunkt (N.\,\ref{60337}). In einem differentialgeometrischen Ansatz bestimmt er die Verteilung der Kräfte beim Zug an gespannten Fäden (N.\,\ref{60117}).%
\\ \indent
Ergänzende Untersuchungen befassen sich mit der Trägheit von Körpern bei schneller Bewegung (N.\,\ref{41172}) und dem Verhalten verbundener Körper unter Krafteinwirkung (N.\,\ref{60277}). Ein allgemeines Beweisprinzip für kontinuierliche Veränderungen wird am Beispiel der Spannungsverteilung in einem aufgehängten Seil erarbeitet (N.\,\ref{39978}).%
\par
In Zusammenhang mit Flüssigkeiten sind die Kapillarphänomene und die Mechanik des Eindringens ein Thema für Leibniz. Er interessiert sich für den Safttransport in Pflanzen, den er sich analog zum Aufsteigen von Flüssigkeiten in engen Röhren erklärt (N.\,\ref{42488}). Das Phänomen der Brechung beim Eindringen fester Körper in Flüssigkeiten ergibt sich laut Leibniz aus unendlich kleinen beweglichen Kügelchen, die diesen Vorgang ermöglichen (N.\,\ref{58952}). %
Eine eigene Betrachtung schenkt Leibniz der Zähigkeit von Flüssigkeiten (N.\,\ref{58218}). Er vergleicht die Bewegungsübertragung durch ein spannbares Seil mit der zwischen Körpern durch Reibung und erstellt ein mechanisches Modell zur Beschreibung der Flüssigkeitsviskosität.%
\par
\vspace{2.0em}
\noindent
II. Statica (N.\,\ref{60066}--N.\,\ref{60128})%
\par
\vspace{0.25em}
\noindent
Die Stücke zur Statik bilden den umfangreichsten Teil der Schriften in diesem Band. Sie zeugen von Leibniz' Auseinandersetzung mit grundlegenden Fragen der Gleichgewichtslehre und seinen Anstrengungen, physikalische Probleme mathematisch zu durchdringen. Im Zentrum steht dabei die Entwicklung strenger geometrischer Beweise für fundamentale Prinzipien der Statik.%
\par
\vspace{1em}
\noindent
II.\;1. Schwerpunkt%
\par\noindent
Ein Hauptaugenmerk liegt auf der theoretischen Fundierung des Schwerpunktbegriffs. In mehreren zusammenhängenden Stücken erarbeitet Leibniz Beweise für die Existenz des Schwerpunkts: Ausgehend von der Annahme eines gemeinsamen Strebepunkts aller Körper (N.\,\ref{60066}) führt er einen geometrischen Beweis, der auf Winkelbetrachtungen und Proportionen basiert (N.\,\ref{60279}). Durch fortgesetzte Halbierung der Winkel zwischen Gleichgewichtsachsen zeigt er, dass jede durch den Schwerpunkt gehende Gerade eine Gleichgewichtsachse sein muss (N.\,\ref{57032}); aus dieser Betrachtung von Gleichgewichtsachsen gewinnt Leibniz eine Beweisführung für die Existenz eines gemeinsamen Schwerpunkts in beliebigen Körpern (N.\,\ref{41137}). %
Des Weiteren befasst er sich mit der Schwerpunktbestimmung bei gekrümmten Rotationsachsen und wählt einen geometrischen Ansatz, um den Schwerpunkt von Kurven bei ihrer Entwicklung zu finden (N.\,\ref{Cc2-1421}). Seine Studien erstrecken sich auch auf die Unterscheidung zwischen relativem und absolutem Schwerpunkt bei Kurven (N.\,\ref{58240}) sowie auf die besondere Problematik der Schwerpunktbestimmung bei Rotationskörpern (N.\,\ref{58657}). In weiteren geometrischen Betrachtungen (N.\,\ref{60309}) zeigt er, dass der Schwerpunkt eines Halbkreises mit dem Schwerpunkt eines konzentrischen Kreisbogens zusammenfällt und überträgt diese Erkenntnis auf Halbkugeln und deren Oberflächen. Eine eigene Gruppe bilden die Arbeiten zum Gleichgewichtsprinzip (N.\,\ref{57033+41135}, N.\,\ref{41134}, N.\,\ref{41136}): Hier entwickelt Leibniz eine Beweisführung, die sich bewusst von der archimedischen Tradition absetzt, wobei er rein geometrisch verfährt. In diesen französischsprachigen Stücken hebt er wiederholt hervor, dass sein analytischer Ansatz ohne weitere Vorannahmen den Weg weist, wie  das richtige Gleichgewichtsverhältnis zu entdecken ist, und sich dadurch gegenüber algebraischen Ansätzen auszeichnet. Die späteren Arbeiten zeigen eine zunehmende Verallgemeinerung der statischen Prinzipien. Bei der Behandlung von Kräften, deren Wirkungslinien sich außerhalb des Schwerpunkts schneiden (N.\,\ref{60314}), beweist Leibniz die universelle Gültigkeit der Proportionen unabhängig vom Schnittpunkt.%
\par
\vspace{1em}
\noindent
II.\;2. Schiefe Ebene%
\par \noindent
Besonders intensiv befasst sich Leibniz mit der Theorie der schiefen Ebene, vor allem in kritischer Auseinandersetzung mit Giovanni Francesco Vanni\indexN{\textls{Vanni}, Giovanni Francesco 1638\textendash1709}, die sich über mehrere Jahre erstreckt. Dessen Einwände gegen die etablierte statische Theorie beantwortet Leibniz mit einer Reihe geometrischer Beweise (N.\,\ref{41667}, N.\,\ref{55606}). Eine zweite Phase seiner Auseinandersetzung mit Vanni führt über mehrere Anläufe (N.\,\ref{32310}, N.\,\ref{60341}, \ref{32313}) hin zu einer für die \textit{Acta Eruditorum} verfassten Darlegung des geometrischen Beweises, mit dem Leibniz vollständig auf die Kritik Vannis\indexN{\textls{Vanni}, Giovanni Francesco 1638\textendash1709} antwortet (N.\,\ref{49409}). Darin löst Leibniz das von Vanni vorgebrachte Problem einer Kugel auf zwei senkrecht stehenden Ebenen, indem er die \textit{Regula alternativorum} einführt und damit die Verteilung der Gesamtkraft auf die beiden Ebenen bestimmt. In späteren Rezensionen zu Werken (N.\,\ref{61037}, N.\,\ref{41824}) von Vanni verfolgt Leibniz dessen Diskussion der Statik weiter. Seine anhaltende Beschäftigung mit dieser Problematik dokumentieren weitere Theoreme zu den Kräfteverhältnissen bei geneigten Geraden (N.\,\ref{60030}, N.\,\ref{60100}). Diese Ausgangsbedingung erweitert Leibniz auf den Fall einer Kugel, die auf zwei bewegliche Schenkel eines Winkels drückt; er bestimmt dafür, welche Gewichte erforderlich sind, damit eine Drehung der Schenkel verhindert wird (N.\,\ref{60110}).%
\\ \indent
Die Problematik der schiefen Ebene kehrt in verschiedenen Kontexten wieder. In seiner Auseinandersetzung mit zwei Abhandlungen von Francesco Spoleti\indexN{\textls{Spoleti}, Francesco, gest. 1712} (N.\,\ref{60458}, N.\,\ref{41200}) ist sein Interesse an medizinischen Phänomenen zu erkennen, die sich durch Anwendung mechanischer Prinzipien erklären lassen. Konkret verbindet Leibniz hier seine geometrische Behandlung der Kräfte auf der schiefen Ebene mit physiologischen Überlegungen zur Gallenabsonderung in der Leber. Für sein Betrachtung der schiefen Ebene bezieht er die Hebelgesetz mit ein, indem er sich das spezielle Problem einer Wage vornimmt, auf die ein Gewicht über einen schiefen Hebelarm wirkt (N.\,\ref{60197}).%
\\ \indent
Die schiefe Ebene zieht Leibniz zur Betrachtung weiterer Fälle heran: Er denkt sich eine Kugel, die nicht rollt bzw. am Rollen gehindert wird und ruht (N.\,\ref{60112}) oder ohne Rotation hinuntergleitet (N.\,\ref{60115}, N.\,\ref{60220}). Hierfür berücksichtigt er verschiedene Positionen des Schwerpunkts zur Ebene und zu den stützenden Kräften, und gewinnt eine mathematische Lösung, um die Kräfteverhältnisse zu bestimmen. Auch das Verhalten eines Würfels untersucht Leibniz auf der schiefen Ebene: Er stellt zehn Bedingungen auf, unter denen der Würfel gleitet oder kippt, und bestimmt dafür einen kritischen Neigungswinkel von 45 Grad (N.\,\ref{60118}). Zu den speziellen Untersuchungen gehört das klassische Problem der Kettenlinie, dessen Betrachtung Leibniz durch zusätzliche Gewichte an den Enden der Kette erweitert (N.\,\ref{60128}).%
\par
\vspace{1.0em}
\noindent
II.\;3. Schwerkraft und Technik%
\par \noindent
Von eigenem Interesse ist Leibniz' Auseinandersetzung mit zeitgenössischen Theorien, die sich in mehreren Rezensionen und kritischen Kommentaren niederschlägt. Seine Rezension zu Vannis\indexN{\textls{Vanni}, Giovanni Francesco 1638\textendash1709} \textit{Exegeses physico-mathematicae} (N.\,\ref{41199}) würdigt dessen Scharfsinn bei der Formulierung von Einwänden, zeigt aber zugleich die Schwächen in dessen Argumentationen auf. In seinen Kommentaren zu Auszügen aus Pierre Varignons\indexN{\textls{Varignon}, Pierre de 1654\textendash1722} \textit{Projet d'une nouvelle mechanique} (N.\,\ref{60365}) befasst er sich mit dessen neuer Methode zur Behandlung mechanischer Probleme, insbesondere die Anwendung des Kräfteparallelogramms zur Bestimmung von Gleichgewichtszuständen. %
Leibniz' Auseinandersetzung mit der wissenschaftlichen Literatur seiner Zeit geht aus weiteren Auszügen und Kommentaren hervor. Aus Pierre Fermats\indexN{\textls{Fermat}, Pierre de 1601\textendash1665} \textit{Varia opera mathematica} (N.\,\ref{58229}) exzerpiert er dessen Behandlung mechanischer Probleme und kritisiert die Aussagen zum Verhältnis zwischen Gewicht und Schwerkraftrichtung. Bei der Lektüre von Vitale Giordanis\indexN{\textls{Giordani} (Giordano), Vitale 1633\textendash1711} Werken (N.\,\ref{58253}) konzentriert er sich auf dessen Beweise zum Verhältnis von Gewichten auf schiefen Ebenen, wobei er mögliche Schwachstellen in der Argumentation identifiziert.%
\\ \indent
Eine kurze Aufzeichnung (N.\,\ref{60114}) gilt zunächst der Schwerkraft auf der schiefen Ebene und zeigt, dass die Wirkung eines schweren Körpers auf die Ebene äquivalent ist zu einer Kraft, die senkrecht zur Ebene wirkt. In einem weiteren Stück (N.\,\ref{60631}) leitet Leibniz mathematisch her, wie sich die Gewichtskraft eines schweren Körpers auf zwei geneigte Ebenen verteilt. Bei der Behandlung von Kräften, die eine schwere Kugel auf zwei stützende Ebenen ausübt (N.\,\ref{60113}), bestimmt er mittels einer geometrischen Konstruktion die Kräfteverhältnisse; dabei zeigt sich, dass die Form des Körpers keine Rolle spielt, solange die Normalen zu den Ebenen durch den Schwerpunkt gehen. Diese Erkenntnis verallgemeinert er in seiner Einschätzung zur Kraft, die nötig ist, um eine Last auf einer Ebene zu halten und zu bewegen (N.\,\ref{60116}), wobei er verschiedene Fälle der Lastverteilung und Bewegung betrachtet.%
\\ \indent
Eine eigene Gruppe von Stücken beschäftigt sich mit praktischer Mechanik, in denen Leibniz technische Anwendungen der statischen Prinzipien untersucht. Aus seinen vergleichenden Studien zu Zahnrädern und verbundenen Hebeln gewinnt er Regeln für eine optimale Kraftübertragung (N.\,\ref{57034}); darin fließen seinen gewonnenen Erkenntnisse über den Eingriffwinkel, die Form und Ausrichtung der Zähne sowie die Vermeidung von Reibung ein. %
Drei Stücke beschäftigen sich mit Seilsystemen und deren Gleichgewichtsbedingungen. Leibniz betrachtet ein Gewicht, das von zwei Seilen getragen wird, berechnet deren Spannung in schiefer Kraftrichtung und bestimmt die optimale Position des Schwerpunkts  (N.\,\ref{58943}). Im weiteren Verlauf berechnet Leibniz für ein Gewicht, das an mehreren Seilen aufgehängt ist, das Verhältnis zwischen dem Gewicht und den Haltekräften der Seile im Gleichgewichtszustand  (N.\,\ref{60105}). Seine Berechnungen zeigen, dass die Seilkräfte von den Winkeln der Seile zur Vertikalen abhängen.%
\par
\vspace{2.0em}
\noindent
III. Stoß (N.\,\ref{60221-B}--N.\,\ref{58950})%
\par
\vspace{0.25em}
\noindent
Die Stücke zur Theorie des Stoßes zeigen Leibniz' anhaltende Auseinandersetzung mit diesem zentralen Problem der Mechanik. Die in diesem Band edierten Stücke sind (ganz überwiegend) ab den 1680er Jahren entstanden. Leibniz' Beschäftigung mit diesem Thema reicht bis auf die Pariser Zeit zurück und ist durch eine Vielzahl von Stücken dokumentiert, die in den Bänden zwei und drei der Reihe VIII erschienen sind.
\par
\vspace{1.0em}
\noindent
III.\;1. Einfacher Stoß%
\par \noindent
Die theoretischen Grundlagen schafft Leibniz in einer Reihe von Stücken, die sich der Bestimmung der relativen Kraft beim Zusammenstoß von Körpern widmen (N.\,\ref{60243}). Dabei überlegt er, wie die Kraft beim Aufprall eines schnelleren auf einen langsameren Körper zu bemessen ist, und kommt zu dem Schluss, dass die relative Geschwindigkeit maßgeblich ist. Eine grundlegende Kritik der cartesischen Mechanik liefert Leibniz in seiner Auseinandersetzung mit den sieben Stoßregeln von Descartes\indexN{\textls{Descartes} (Cartesius, des Cartes), René 1596\textendash1650} (N.\,\ref{60099}), deren Widersprüche er durch Gegenbeispiele aufzeigt. %
Indem er der Frage nachgeht, wie sich beim Zusammentreffen zweier Körper die relative Kraft in eine Kreisbewegung umwandelt, gelangt er zu einem genaueren Verständnis der Bewegungsübertragung (N.\,\ref{60111}). Parallel dazu formuliert Leibniz das fundamentale Prinzip der Richtungserhaltung (N.\,\ref{58220}) und differenziert zwischen absoluter und beharrender Kraft (\textit{potentia conservativa}) beim Zusammenstoß (N.\,\ref{58235}). %
Aus seinen Studien zum Stoß gewinnt er weitere Konzepte für dessen differenzierte Betrachtung: Die Kontaktlinie zwischen den stoßenden Körpern und deren Stoßrichtung, auf die er die resultierende Bewegung zurückführt (N.\,\ref{60356}).%
\par
\vspace{1.0em}
\noindent
III.\;2. Exzentrischer Stoß%
\par \noindent
Ein besonderer Fall ist der exzentrische Stoß, mit dem sich Leibniz mehrfach befasst. In zwei Stücken (N.\,\ref{60142+60297}, N.\,\ref{60299}) geht es um die Wirkung eines Stoßes, dessen Richtung nicht durch den Schwerpunkt geht; Leibniz trifft hieraus geometrisch eine Aussage über die resultierende Kreisbewegung. Diese Betrachtung vertieft er in zwei eigenhändig datierten und gezählten Bogen (N.\,\ref{58940}, N.\,\ref{58939}), wobei er besonders auf die Rolle des Schwerpunkts eingeht. Das Geheimnis des exzentrischen Stoßes will Leibniz nach eigenen Worten dadurch gelüftet haben, dass er die gestoßenen Körper gedanklich in Teile zerlegt, um die in Translation und Rotation aufgeteilte Bewegung genauer zu bestimmen (N.\,\ref{60339}).
Seine Analyse des exzentrischen Stoßes vertieft und differenziert er davon ausgehend weiter: Bei der resultierenden Kombination aus Vorwärts- und Drehbewegung berücksichtigt er die Rolle des Schwerpunkts und dessen Abstand zum Stoßpunkt (N.\,\ref{60300}, N.\,\ref{60106}, N.\,\ref{60107}).
\par
\vspace{1.0em}
\noindent
III.\;3. Mehrkörperstoß%
\par \noindent
Die Komplexität des Phänomens Stoß steigert sich in der Untersuchung von Fällen, in denen mehr als zwei Körper aufeinandertreffen. Die Bewegung des gemeinsamen Schwerpunkts (N.\,\ref{66097}) und die Zusammensetzung der Bewegungen beim Stoß mehrerer Körper (N.\,\ref{60295}) bilden dabei zentrale Themen. Leibniz formuliert das Prinzip der Erhaltung der relativen Kraft zwischen drei Körpern (N.\,\ref{53450}) und liefert eine neue Methode zur Analyse des Dreikörperstoßes (N.\,\ref{60298}). Die Frage nach der Erhaltung der relativen Kraft beim schiefen Stoß führt zu weiteren methodischen Überlegungen (N.\,\ref{58951}).%
\par
\vspace{1.0em}
\noindent
III.\;4. Allgemeine Stoßtheorie%
\par \noindent
Seine späteren Arbeiten widmen sich verstärkt methodologischen Fragen. Er entwickelt einen allgemeinen Ansatz für Phänomene des Stoßes (N.\,\ref{58947}). Die mathematische Behandlung des Problems verfeinert er in mehreren Stücken (N.\,\ref{58300}, N.\,\ref{60359}), auch für den Fall des Mehrkörperstoßes (N.\,\ref{58293}, N.\,\ref{39763}). Er betrachtet die Geschwindigkeitsänderungen beim Stoß (N.\,\ref{60358}, N.\,\ref{60362}) und die Bedingungen für eine Bewegungsübertragung, zu denen er die Elastizität und das Verhältnis der Annäherungs- und Rückprallgeschwindigkeiten zählt (N.\,\ref{60360}).%
\\ \indent
Die Frage nach den Grundprinzipien der Stoßtheorie beschäftigt ihn in mehreren Stücken. Er prüft die Rolle der Krafterhaltung (N.\,\ref{41780}) und die Möglichkeit, alle Stoßphänomene aus wenigen Grundsätzen abzuleiten. Seine methodischen Reflexionen münden in eine Demonstration zweier Hypothesen zur Erklärung der Bewegungsgesetze (N.\,\ref{60354}). Den Abschluss bildet eine umfassende Analyse des Mehrkörperstoßes (N.\,\ref{42014}), die seine theoretischen Einsichten mit praktischen Anwendungen verbindet.
\\ \indent
Leibniz ist es in mehreren Stücken ein Anliegen, die Phänomene des Stoßes zu differenzieren, nach verschiedenen Arten zu klassifizieren und die spezifischen Eigenschaften zu erfassen (N.\,\ref{41121}, N.\,\ref{58946}, N.\,\ref{58950}). Seine Methode zur Bestimmung der relativen Kraft (N.\,\ref{58949}) und seine Beschäftigung mit den Grundgesetzen des Stoßes ergänzen diese systematische Behandlung des Themas. Die Zeichnungen zum zeitgenössichen Billardspiel veranschaulichen verschiedene Stoßkonstellationen, bei denen eine Kugel über Bande oder vermittelst einfachen oder mehrfachen Stoß durch das aufgestellte Zieltor gehen kann (N.\,\ref{60221-B}).  
\par
\vspace*{3em}
\hspace{104mm}Harald Siebert
